% PROOFS
\section{Mathematical Formulations}%
\label{sec:mathematical_formulations}

To generate the transition probabilities for the Markov chain usage model, the symbolic layer of NeSy-MBST models the state-transition structure as a constrained optimization problem~\cite{whittaker1993markov, boehr_constraints}. This section formalises the mathematical foundations underlying the optimizer.

Let $\mathcal{S} = \{s_1, s_2, \dots, s_n\}$ be the finite set of usage states identified by the $L^{*}$ active learning process~\cite{lstar_llm}. The stochastic matrix $\mathbf{P} = [p_{ij}] \in \mathbb{R}^{n \times n}$ represents the transition probabilities, where $p_{ij}$ denotes the probability of transitioning from state~$s_i$ to state~$s_j$. The optimizer must compute a feasible~$\mathbf{P}$ that simultaneously satisfies several families of structural and linear constraints, detailed in the following subsections.

% Subsection A: Primary Probability Axioms
\subsection{Primary Probability Axioms}%
\label{subsec:probability_axioms}

Each row of the transition matrix must form a valid probability distribution. Formally, two axioms must hold:

\begin{equation}
\label{eq:nonnegativity}
0 \leq p_{ij} \leq 1, \quad \forall\; i,j \in \{1, 2, \dots, n\},
\end{equation}

\begin{equation}
\label{eq:row_stochastic}
\sum_{j=1}^{n} p_{ij} = 1, \quad \forall\; i \in \{1, 2, \dots, n\}.
\end{equation}

Constraint~\eqref{eq:nonnegativity} enforces non-negativity and boundedness of each entry, while constraint~\eqref{eq:row_stochastic} ensures that $\mathbf{P}$ is row-stochastic, i.e., the outgoing probabilities from every state sum to unity~\cite{nap_statistics}.

% Subsection B: Structural Absence Constraints
\subsection{Structural Absence Constraints}%
\label{subsec:structural_constraints}

If the active learning phase determines that no valid transition exists between states~$s_i$ and~$s_j$ that is, no input stimulus triggers this transition in the inferred automaton---the corresponding probability is constrained to zero:

\begin{equation}
\label{eq:structural_zero}
p_{ij} = 0, \quad \forall\; (i,j) \notin \mathcal{E},
\end{equation}

\noindent where $\mathcal{E} \subseteq \mathcal{S} \times \mathcal{S}$ represents the set of verified transition edges in the state graph. These structural absence constraints ensure that the resulting Markov chain is topologically consistent with the automaton learned during the $L^{*}$ inference step~\cite{lstar_llm}. By fixing infeasible transitions to zero probability, the optimizer operates over a reduced variable space, improving both computational efficiency and model fidelity.

% Subsection C: Evolving Operational Constraints
\subsection{Evolving Operational Constraints}%
\label{subsec:operational_constraints}

Comparative constraints extracted by the LLM from natural-language requirements are formulated as linear relationships between transition probabilities~\cite{boehr_constraints}:

\begin{equation}
\label{eq:proportional}
p_{ij} = \alpha \cdot p_{kl},
\end{equation}

\noindent where $\alpha \in \mathbb{R}_{>0}$ is a positive scaling factor representing the proportional likelihood of one user action relative to another. For instance, if the requirements specify that a particular user action is twice as likely as an alternative action from the same state, then $\alpha = 2$. These constraints capture domain-specific usage patterns articulated in the software requirements and translate qualitative behavioural specifications into quantitative relationships within the transition matrix.

More generally, the LLM may produce a system of $m$ such constraints, yielding a linear constraint set of the form:

\begin{equation}
\label{eq:linear_system}
\mathbf{A}\,\mathrm{vec}(\mathbf{P}) = \mathbf{b}, \quad \mathbf{A} \in \mathbb{R}^{m \times n^2},\; \mathbf{b} \in \mathbb{R}^{m},
\end{equation}

\noindent where $\mathrm{vec}(\mathbf{P})$ denotes the vectorisation of~$\mathbf{P}$ and each row of $\mathbf{A}$ encodes one proportional or equality constraint extracted from the requirements~\cite{emml_usage_models}.

% Subsection D: Long-Run and Expected Value Constraints
\subsection{Long-Run and Expected Value Constraints}%
\label{subsec:longrun_constraints}

Beyond instantaneous transition probabilities, the usage model must also respect constraints on long-run behaviour and expected traversal costs~\cite{markov_reliability}.

\subsubsection{Steady-State Occupancy}
The steady-state (stationary) distribution vector $\boldsymbol{\pi} = [\pi_1, \pi_2, \dots, \pi_n]$ satisfies:

\begin{equation}
\label{eq:steady_state}
\boldsymbol{\pi}\,\mathbf{P} = \boldsymbol{\pi}, \quad \sum_{i=1}^{n} \pi_i = 1, \quad \pi_i \geq 0 \;\;\forall\; i.
\end{equation}

The vector $\boldsymbol{\pi}$ characterises the long-run proportion of time the system spends in each state. Convex constraints on the occupancy vector take the form:

\begin{equation}
\label{eq:occupancy_bounds}
L_i \leq \pi_i \leq U_i, \quad \forall\; i \in \{1, 2, \dots, n\},
\end{equation}

\noindent where $L_i$ and $U_i$ represent the lower and upper bounds, respectively, on the expected occupancy of state~$s_i$. These bounds are derived from operational profiles or domain expertise and ensure that the generated test suite allocates realistic proportions of effort to each functional area~\cite{whittaker1993markov,whittaker1994markov}.

\subsubsection{Mean First Passage Times}
The expected number of transitions (mean first passage time) from state~$s_i$ to state~$s_j$, denoted $m_{ij}$, is defined recursively as:

\begin{equation}
\label{eq:first_passage}
m_{ij} = 1 + \sum_{\substack{k=1 \\ k \neq j}}^{n} p_{ik} \cdot m_{kj}, \quad \forall\; i \neq j.
\end{equation}

To bound the expected traversal cost ensuring that test sequences remain tractable in length an upper-bound constraint is imposed:

\begin{equation}
\label{eq:cost_bound}
m_{ij} \leq T_{\max},
\end{equation}

\noindent where $T_{\max}$ is a user-specified maximum on the expected number of steps between designated states. This constraint prevents pathologically long expected test paths and ensures practical executability of the resulting test suites~\cite{nap_statistics}.

% Optimisation Objective
\subsection{Entropy-Maximising Objective}%
\label{subsec:objective}

Given the constraint families defined in Sections~\ref{subsec:probability_axioms}--\ref{subsec:longrun_constraints}, the system employs convex programming to compute a feasible transition matrix~$\mathbf{P}$. To avoid unintended testing bias when the constraints admit multiple feasible solutions, the optimizer maximises the entropy of the transition matrix:

\begin{equation}
\label{eq:entropy_objective}
\max_{\mathbf{P}} \;\; H(\mathbf{P}) = -\sum_{i=1}^{n} \sum_{j=1}^{n} p_{ij} \ln p_{ij},
\end{equation}

\noindent subject to constraints~\eqref{eq:nonnegativity}--\eqref{eq:cost_bound}. Maximising entropy yields the least-biased distribution consistent with all known constraints, following the principle of maximum entropy~\cite{emml_usage_models}. The resulting matrix~$\mathbf{P}$ is then used directly to parameterise the Markov chain usage model for statistical test-case generation.
